\setcounter{section}{4}
\setcounter{subsection}{0}

\section{Teoremas do Anulamento, Valor Intermediário e Weierstrass}

\begin{exercicio}[1]
Seja $f(x) = x^5 + x + 1$. Justifique a afirmação: $f$ tem pelo menos uma raiz no intervalo $[-1, 0]$.
\end{exercicio}

\begin{exercicio}[2]
Prove que a equação $x^3 - 4x + 2 = 0$ admite três raízes reais distintas.
\end{exercicio}

\begin{exercicio}[3]
Seja $\alpha$ a menor raiz positiva da equação $x^3 - 4x + 2 = 0$. Determine intervalos de amplitudes $\dfrac{1}{2}$, $\dfrac{1}{4}$ e $\dfrac{1}{8}$ que contenham $\alpha$.
\end{exercicio}

\begin{exercicio}[4]
Prove que a equação $x^3 - \dfrac{1}{1 + x^4} = 0$ admite ao menos uma raiz real.
\end{exercicio}

\begin{exercicio}[5]
Prove que cada um dos conjuntos abaixo admite máximo e mínimo:
\begin{enumerate}[label=\alph*)]
    \item $A = \left\{ \dfrac{x}{1 + x^2} \;\middle|\; -2 \le x \le 2 \right\}$
    \item $A = \left\{ \dfrac{x^2 + x}{1 + x^2} \;\middle|\; -1 \le x \le 1 \right\}$
\end{enumerate}
\end{exercicio}

\begin{exercicio}[6]
Seja $f\colon [-1, 1] \to \R$ dada por $f(x) = \dfrac{x^2 + x}{1 + x^2}$.
\begin{enumerate}[label=\alph*)]
    \item Prove que $f(1)$ é o valor máximo de $f$.
    \item Prove que existe $x_1 \in {]-1, 0[}$ tal que $f(x_1)$ é o valor mínimo de $f$.
\end{enumerate}
\end{exercicio}

\begin{exercicio}[7]
\begin{enumerate}[label=\alph*)]
    \item Prove que todo polinômio de grau 3 admite pelo menos uma raiz real.
    \item Prove que todo polinômio de grau ímpar admite pelo menos uma raiz real.
\end{enumerate}
\end{exercicio}

\begin{exercicio}[8]
Seja $f\colon [a, b] \to \R$ uma função contínua e suponha que $f$ não seja constante em $[a, b]$. Prove que existem números reais $m$ e $M$, com $m < M$, tais que $\operatorname{Im} f = [m, M]$.

\smallskip
\noindent\textbf{Observação:} Imagem de $f = \operatorname{Im} f = \{f(x) \mid x \in [a, b]\}$.
\end{exercicio}

\begin{exercicio}[9]
Seja $f\colon I \to \R$ contínua, em que $I$ é um intervalo qualquer. Prove que a imagem de $f$ é um intervalo.
\end{exercicio}

\begin{exercicio}[10]
Suponha que $f\colon [0, 1] \to \R$ seja contínua, $f(0) = 1$ e que $f(x)$ é racional para todo $x \in [0, 1]$. Prove que $f(x) = 1$, para todo $x \in [0, 1]$.
\end{exercicio}

\begin{exercicio}[11]
Seja $f\colon [0, 1] \to \R$ contínua e tal que, para todo $x \in [0, 1]$, $0 \le f(x) \le 1$. Prove que existe $c \in [0, 1]$ tal que $f(c) = c$.
\end{exercicio}

\begin{exercicio}[12]
Seja $f$ contínua em $[a, b]$ e tal que $f(a) < f(b)$. Suponha que quaisquer que sejam $s$ e $t$ em $[a, b]$, $s \neq t \implies f(s) \neq f(t)$. Prove que $f$ é estritamente crescente em $[a, b]$.

\smallskip
\noindent\textbf{Observação:} $f$ estritamente crescente em $[a, b] \iff \forall s, t \in [a, b],\; s < t \implies f(s) < f(t)$.
\end{exercicio}

\begin{exercicio}[13]
Suponha $f$ contínua no intervalo $I$ e que $f$ admita neste intervalo uma única raiz $a$. Suponha, ainda, que existe $x_0 \in I$, com $x_0 > a$, tal que $f(x_0) > 0$. Prove que, para todo $x \in I$, com $x > a$, $f(x) > 0$.
\end{exercicio}

\begin{exercicio}[14]
Considere a função $f$ dada por:
\[
f(x) = 2x^3 - \sqrt{x^2 + 3x}.
\]
\begin{enumerate}[label=\alph*)]
    \item Verifique que $f$ é contínua em $[0, +\infty[$.
    \item Mostre que $1$ é a única raiz de $f$ em ${]0, +\infty[}$, que $f(2) > 0$ e que $f\!\left(\dfrac{1}{2}\right) < 0$.
    \item Conclua que $f(x) > 0$ em ${]1, +\infty[}$ e que $f(x) < 0$ em ${]0, 1[}$.
\end{enumerate}
\end{exercicio}

\begin{exercicio}[15]
Suponha $f$ contínua em $I$ e sejam $a$ e $b$ pertencentes a $I$, com $a < b$, as únicas raízes de $f$ em $I$. Sejam $x_0, x_1 \text{ e } x_2$ em $I$ com $x_0 < a$, $a < x_1 < b$ e $b < x_2$. Estude o sinal de $f$ em $I$, a partir dos sinais de $f(x_0)$, $f(x_1)$ e $f(x_2)$. Justifique.
\end{exercicio}
